\subsubsection{Categories of modeler-supplied context}
\label{sec:curation-categories}

The field-level curation rates above aggregate changes of several different kinds. Across the SubmodelTarget and CalibrationTarget revisions, modeler-supplied context fell into the following recurring categories:

\begin{enumerate}
  \item \emph{Unit conventions and dimensional bridges.} Conversion factors that the LLM cannot infer from \texttt{model\_structure.json} alone, such as a missing $10^{-6}$ tissue protein density conversion in cDC and Treg density formulas, or the Abercrombie 3D-to-2D correction used to convert section-level cell counts to volumetric densities.
  \item \emph{Source selection and retirement.} Excluding a study after recognizing a mismatch between assay or marker definitions and model categorization (e.g., a study reporting CD86+ as an M1 marker, where the simulation does not enforce that definition), or substituting a better source when one becomes available.
  \item \emph{Cross-source aggregation.} Decisions to pool measurements from multiple studies (e.g., \texttt{cd8\_density} aggregated across Golesworthy 2022 and Liu 2015) rather than treat them as separate targets.
  \item \emph{Measurement-to-state translation logic.} Code inside the observable function that converts a paper's measurement into a model-relevant quantity, such as serum-to-tissue bridging factors that translate plasma cytokine concentrations to tissue-level model species, or active-fraction adjustments where immunoassays detect total cytokine although only a signalling-active fraction is biologically meaningful.
  \item \emph{Denominator and normalization choices.} What a percentage is taken of (CD8 of CD3 versus of nucleated cells), and which compartment volume normalizes a per-tissue measurement.
  \item \emph{Scenario-specific forward-model structure.} Constraints particular to a calibration scenario rather than to the underlying biology: peak-ordering constraints (e.g., a tumor-volume minimum that must follow a CD8 peak in the GVAX-only arm), or reparameterization of a kinetic in its linear sub-saturating regime when the data do not span saturation.
  \item \emph{Uncertainty widening from source quality.} Adjusting prior or likelihood widths upward when a source is cross-species or cross-indication; the schema's \texttt{source\_relevance} block records the rationale alongside the widening.
\end{enumerate}

These categories did not contribute equally to curation effort. Unit conventions and source selection accounted for the largest share of edits. Auxiliary-parameter additions and denominator decisions occurred less often but represented modeling decisions about how a measurement maps to a model state, and typically took longer per edit than the unit corrections. Forward-model restructuring and uncertainty widening were applied selectively to a smaller number of targets.
